\chapter{Simetri Pusat dan Jajargenjang}\label{ch:g7:central}

Putarlah sebuah bangun setengah putaran mengelilingi sebuah titik:
inilah \emph{\omterm{def:g7:central:def}{simetri pusat}}, transformasi kedua dalam buku ini sesudah
pencerminan pada \cref{ch:g6:symmetry}. Bangun bintangnya adalah
\omterm{def:g7:central:parallelogram}{jajargenjang} --- \omterm{def:g4:geometry:polygon}{segi empat} yang \omterm{def:g7:central:def}{simetris} terhadap titik potong
diagonalnya sendiri.

\section{Simetri pusat}

\begin{definition}[Simetris terhadap sebuah titik]\label{def:g7:central:def}
Titik \emph{simetris}\index{simetri pusat} dari titik $M$ terhadap
titik $O$ adalah titik $M'$ sehingga $O$ menjadi \emph{titik tengah}
$[MM']$. Bangun yang simetris adalah bangun asalnya sesudah setengah
putaran ($180^\circ$) mengelilingi $O$.
\end{definition}

\begin{omfigure}
\begin{tikzpicture}[scale=1.0]
  \fill (0,0) circle (1.6pt) node[below, font=\small] {$O$};
  \fill[omDef] (-2.3,1.0) circle (1.8pt) node[above, font=\small] {$M$};
  \fill[omThm] (2.3,-1.0) circle (1.8pt) node[below, font=\small] {$M'$};
  \draw[dashed] (-2.3,1.0) -- (2.3,-1.0);
  \draw (-1.2,0.45) -- (-1.1,0.6);
  \draw (1.1,-0.6) -- (1.2,-0.45);
  \fill[omDef] (-1.4,-0.8) circle (1.8pt) node[below, font=\small] {$N$};
  \fill[omThm] (1.4,0.8) circle (1.8pt) node[above, font=\small] {$N'$};
  \draw[dashed] (-1.4,-0.8) -- (1.4,0.8);
\end{tikzpicture}

{\small Simetri terhadap $O$: setiap titik dan bayangannya segaris
dengan $O$, berjarak sama di kedua sisinya (tanda garis).}
\end{omfigure}

\begin{method}[Melukis titik simetris]\label{met:g7:central:construct}
Untuk melukis titik \omterm{def:g7:central:def}{simetris} $M'$ dari $M$ terhadap $O$:
\begin{enumerate}
  \item gambarlah garis $(MO)$ lalu perpanjanglah melewati $O$;
  \item ukurlah $MO$ (dengan penggaris atau jangka);
  \item letakkan $M'$ pada perpanjangannya dengan $OM' = OM$.
\end{enumerate}
Pada kertas berpetak: hitunglah langkah mendatar dan tegak dari $M$ ke
$O$, lalu ulangi langkah yang \emph{sama} dari $O$ untuk mencapai $M'$.
\end{method}

\begin{example}[Pada petak]\label{ex:g7:central:grid}
Kalau $M$ berada $3$ petak ke kiri dan $1$ petak ke atas dari $O$,
titik simetrisnya $M'$ berada $3$ petak ke \emph{kanan} dan $1$ petak
ke \emph{bawah} dari $O$: \omterm{def:g7:central:def}{simetri pusat} membalikkan kedua arahnya
sekaligus (pencerminan hanya membalikkan satu arah).
\end{example}

\begin{proposition}[Apa yang dipertahankan setengah putaran]\label{prop:g7:central:preserved}
\omterm{def:g7:central:def}{Simetri pusat} mempertahankan panjang, sudut, \omterm{def:g6:measure:perimeter}{keliling} dan luas;
bayangan sebuah garis adalah garis yang \emph{\omterm{def:g6:lines:perp}{sejajar}}, bayangan \omterm{def:g6:lines:objects}{ruas
garis} adalah ruas \omterm{def:g6:lines:perp}{garis sejajar} yang sama panjangnya, bayangan
lingkaran adalah lingkaran yang \omterm{def:g3:shapes:shapes}{jari-jarinya} sama (berpusat di titik
\omterm{def:g7:central:def}{simetris} dari pusat lingkarannya).
\end{proposition}
\admitted

\begin{remark}
Bandingkan dengan pencerminan (\cref{prop:g6:symmetry:preserved}):
keduanya mempertahankan panjang dan sudut; tetapi pencerminan
membalikkan bangunnya (huruf ``b'' menjadi ``d''), sedangkan
setengah putaran menjaganya tetap terbaca --- terbalik ke bawah
(huruf ``b'' menjadi ``q''). Dan tidak seperti pencerminan, setengah putaran
mengirim setiap garis ke garis yang \omterm{def:g6:lines:perp}{sejajar} dengannya.
\end{remark}

\begin{definition}[Pusat simetri]\label{def:g7:central:center}
Sebuah titik $O$ disebut \emph{pusat simetri}\index{pusat simetri}
sebuah bangun jika setengah putaran mengelilingi $O$ mengirim
bangunnya tepat ke dirinya sendiri. Misalnya, huruf S, N, Z punya
pusat simetri; lingkaran punya satu (yaitu pusatnya); segitiga tidak
pernah punya.
\end{definition}

\section{Jajargenjang}

\begin{definition}[Jajargenjang]\label{def:g7:central:parallelogram}
\emph{Jajargenjang}\index{jajargenjang} adalah \omterm{def:g4:geometry:polygon}{segi empat} yang
diagonalnya berpotongan di titik tengah bersamanya. Titik potong itu
lalu menjadi \omterm{def:g7:central:center}{pusat simetri} bangunnya: setiap \omterm{def:g5:solids:def}{titik sudutnya} adalah
bayangan setengah putaran dari \omterm{def:g5:solids:def}{titik sudut} yang berseberangan.
\end{definition}

\begin{theorem}[Sifat jajargenjang]\label{thm:g7:central:properties}
Pada \omterm{def:g7:central:parallelogram}{jajargenjang} $ABCD$ dengan pusat $O$:
\begin{enumerate}
  \item sisi yang berhadapan \omterm{def:g6:lines:perp}{sejajar}: $(AB) \parallel (CD)$ dan
        $(BC) \parallel (AD)$;
  \item sisi yang berhadapan sama panjang: $AB = CD$ dan
        $BC = AD$;
  \item sudut yang berhadapan sama besar: $\widehat A = \widehat C$ dan
        $\widehat B = \widehat D$.
\end{enumerate}
\end{theorem}

\begin{proof}
Setengah putaran mengelilingi $O$ mengirim $A$ ke $C$ dan $B$ ke $D$
(menurut definisinya: $O$ adalah titik tengah kedua diagonalnya). Jadi
ia mengirim \omterm{def:g6:lines:objects}{ruas garis} $[AB]$ ke \omterm{def:g6:lines:objects}{ruas garis} $[CD]$; menurut
\cref{prop:g7:central:preserved} bayangannya \omterm{def:g6:lines:perp}{sejajar} dengan $[AB]$ dan
sama panjang: itulah bagian 1 dan 2. Ia juga mengirim sudut di $A$ ke
sudut di $C$, dan sudut dipertahankan: itulah bagian 3.
\end{proof}

\begin{omfigure}
\begin{tikzpicture}[scale=1.0]
  \coordinate (A) at (0,0);
  \coordinate (B) at (3.4,0);
  \coordinate (C) at (4.6,1.9);
  \coordinate (D) at (1.2,1.9);
  \draw[thick] (A) -- (B) -- (C) -- (D) -- cycle;
  \draw[omThm] (A) -- (C);
  \draw[omThm] (B) -- (D);
  \fill (2.3,0.95) circle (1.6pt) node[right=3pt, font=\small] {$O$};
  \node[below left, font=\small] at (A) {$A$};
  \node[below right, font=\small] at (B) {$B$};
  \node[above right, font=\small] at (C) {$C$};
  \node[above left, font=\small] at (D) {$D$};
  \draw (1.62,-0.09) -- (1.78,0.09);
  \draw (2.82,1.81) -- (2.98,1.99);
  \draw (0.53,1.02) -- (0.67,0.88);
  \draw (3.93,1.02) -- (4.07,0.88);
\end{tikzpicture}

{\small Sebuah \omterm{def:g7:central:parallelogram}{jajargenjang} dan pusatnya $O$: diagonalnya (merah)
saling memotong di titik tengahnya, sisi yang berhadapan \omterm{def:g6:lines:perp}{sejajar} dan
sama panjang (tanda garis).}
\end{omfigure}

\begin{theorem}[Mengenali jajargenjang]\label{thm:g7:central:recognize}
Sebuah \omterm{def:g4:geometry:polygon}{segi empat} (yang sisinya tidak saling memotong) adalah
\omterm{def:g7:central:parallelogram}{jajargenjang} begitu \emph{salah satu} keadaan berikut berlaku:
\begin{enumerate}
  \item diagonalnya punya titik tengah yang sama (definisinya);
  \item sisi yang berhadapan \omterm{def:g6:lines:perp}{sejajar} dua-dua;
  \item dua sisi yang berhadapan \omterm{def:g6:lines:perp}{sejajar} \emph{dan} sama panjang.
\end{enumerate}
\end{theorem}
\admitted

\begin{example}\label{ex:g7:central:recognize}
$ABCD$ punya $(AB) \parallel (CD)$ dan $AB = CD = 5$ cm. Kriteria 3
berlaku: $ABCD$ adalah \omterm{def:g7:central:parallelogram}{jajargenjang} --- jadi tanpa pengukuran lain
kita juga tahu bahwa $AD = BC$, $(AD) \parallel (BC)$, dan bahwa
diagonalnya saling memotong di tengahnya.
\end{example}

\begin{remark}[Jajargenjang istimewa]
Persegi panjang, \omterm{def:g6:shapes:quadrilaterals}{belah ketupat} dan persegi pada \cref{ch:g6:shapes}
persis \omterm{def:g7:central:parallelogram}{jajargenjang} dengan satu sifat tambahan: diagonal yang sama
panjang (persegi panjang), diagonal yang \omterm{def:g6:lines:perp}{tegak lurus} (\omterm{def:g6:shapes:quadrilaterals}{belah ketupat}),
keduanya (persegi). Alasan sudut yang dijanjikan pada
\cref{ex:g6:shapes:recognize} kini berjalan: pada \omterm{def:g7:central:parallelogram}{jajargenjang}, sudut
yang berurutan berpelurus (\omterm{def:g7:triangles:pairs}{sudut berseberangan} atau sehadap dengan
\omterm{def:g6:lines:perp}{garis sejajarnya}, \cref{thm:g7:triangles:equalities}), jadi satu \omterm{def:g3:shapes:rightangle}{sudut
siku-siku} memaksa keempatnya.
\end{remark}

\section{Latihan}

\begin{exercise}[$\star$]\label{exo:g7:central:1}
Pada kertas berpetak, letakkan $O$, lalu $M$ empat petak ke kanan dari
$O$ dan satu petak ke atas. Lukislah titik \omterm{def:g7:central:def}{simetris} $M'$ dari $M$
terhadap $O$, dan titik \omterm{def:g7:central:def}{simetris} $M$ terhadap \emph{garis petak tegak}
yang melalui $O$. Bandingkan kedua bayangannya.
\end{exercise}

\begin{exercise}[$\star$]\label{exo:g7:central:2}
Gambarlah segitiga $ABC$ dan sebuah titik $O$ di luarnya. Lukislah
bayangannya $A'B'C'$ oleh setengah putaran mengelilingi $O$. Apa yang
dapat kamu katakan tentang panjang $A'B'$ dan $AB$? Tentang garis
$(A'B')$ dan $(AB)$?
\end{exercise}

\begin{exercise}[$\star$]\label{exo:g7:central:3}
Angka mana, yang ditulis seperti pada layar kalkulator ($0$ sampai
$9$), yang punya \omterm{def:g7:central:center}{pusat simetri}? Huruf besar mana di antara H, A, N, S,
T, Z, O?
\end{exercise}

\begin{exercise}[$\star$]\label{exo:g7:central:4}
$KLMN$ adalah \omterm{def:g7:central:parallelogram}{jajargenjang} berpusat $O$ dengan $KL = 7$ cm,
$LM = 4$ cm dan $\widehat K = 115^\circ$. Sebutkan, dengan alasannya:
$MN$, $NK$, $\widehat M$, dan titik tengah $[LN]$.
\end{exercise}

\begin{exercise}[$\star$]\label{exo:g7:central:5}
Lukislah \omterm{def:g7:central:parallelogram}{jajargenjang} $ABCD$ dari diagonalnya: $AC = 8$ cm dan
$BD = 5$ cm, yang berpotongan di titik tengahnya dengan sudut
$60^\circ$ di antaranya. (Gambarlah diagonalnya lebih dahulu!)
\end{exercise}

\begin{exercise}[$\star$]\label{exo:g7:central:6}
Pada sebuah \omterm{def:g7:central:parallelogram}{jajargenjang}, satu sudutnya berukuran $115^\circ$. Dengan
memakai catatan tentang sudut berurutan yang berpelurus, sebutkan
ketiga sudut lainnya.
\end{exercise}

\begin{exercise}[$\star\star$]\label{exo:g7:central:7}
Gambarlah $A(1, 1)$, $B(4, 2)$ dan titik $O(2.5, 2.5)$. Hitunglah
koordinat bayangan $A'$ dan $B'$ oleh setengah putaran mengelilingi
$O$ (pakailah cara menghitung petak pada \cref{ex:g7:central:grid}).
\end{exercise}

\begin{exercise}[$\star\star$]\label{exo:g7:central:8}
$ABCD$ adalah \omterm{def:g4:geometry:polygon}{segi empat} dengan $AB = CD$. Apakah itu sudah cukup
untuk menjadikannya \omterm{def:g7:central:parallelogram}{jajargenjang}? Kalau tidak, sketsalah contoh
penyangkalnya (trapesium sama kaki membantu), lalu nyatakan apa yang
harus ditambahkan pada hipotesisnya.
\end{exercise}

\begin{exercise}[$\star\star$]\label{exo:g7:central:9}
Misalkan $ABC$ sebuah segitiga dan $O$ titik tengah $[BC]$. Lukislah
titik \omterm{def:g7:central:def}{simetris} $A'$ dari $A$ terhadap $O$. Tunjukkan bahwa $ABA'C$
adalah \omterm{def:g7:central:parallelogram}{jajargenjang}. (Kriteria mana pada
\cref{thm:g7:central:recognize} yang gratis di sini?)
\end{exercise}

\begin{exercise}[$\star\star$]\label{exo:g7:central:10}
Benar atau salah, dengan alasan: ``\omterm{def:g4:geometry:polygon}{segi empat} yang sudut
berhadapannya sama besar dua-dua adalah \omterm{def:g7:central:parallelogram}{jajargenjang}''; ``\omterm{def:g7:central:parallelogram}{jajargenjang}
yang diagonalnya \omterm{def:g6:lines:perp}{tegak lurus} adalah \omterm{def:g6:shapes:quadrilaterals}{belah ketupat}''; ``\omterm{def:g7:central:parallelogram}{jajargenjang}
dengan satu \omterm{def:g3:shapes:rightangle}{sudut siku-siku} adalah persegi panjang''.
\end{exercise}

\begin{exercise}[$\star\star\star$]\label{exo:g7:central:11}
Misalkan $ABCD$ \omterm{def:g4:geometry:polygon}{segi empat} mana pun dan $I$, $J$, $K$, $L$ titik
tengah sisinya $[AB]$, $[BC]$, $[CD]$, $[DA]$. Gambarlah beberapa
contoh. Bangun apa yang selalu tampaknya menjadi $IJKL$? (Buktinya,
dengan teorema titik tengah, menyusul pada \cref{ch:g8:midpoints} ---
dan sekali lagi dengan vektor, di buku Sekolah Menengah.)
\end{exercise}

\section{Soal: Kue, garis berat, dan pusat kedua yang
mustahil}

\begin{problem}[{Soal akhir pekan --- setengah putaran yang bekerja:
potongan kue yang adil, batas untuk garis berat, dan mengapa tidak ada
bangun yang punya dua pusat simetri}]
\label{pb:g7:central:1}
Setengah putaran tampak seperti transformasi yang paling sederhana
--- namun ia memotong kue menjadi dua bagian yang terbukti adil,
mengukur garis berat sebuah segitiga, dan menyembunyikan mutiara kecil
matematika murni: bangun yang terbatas dapat punya satu \omterm{def:g7:central:center}{pusat simetri},
tetapi \emph{tidak pernah dua}. Buktinya memakai penemuan yang akan
kamu buat pada Bagian I: dua setengah putaran, yang dilakukan
berturut-turut, berjumlah menjadi satu geseran.

\medskip
\noindent\textbf{Bagian I --- Senam setengah putaran.}
\begin{enumerate}
  \item Pada kertas berpetak, ambillah titik asal $O(0, 0)$ sebagai
        pusatnya. Dengan cara menghitung petak pada
        \cref{ex:g7:central:grid}, sebutkan bayangan $M(3, 1)$,
        $N(-2, 4)$ dan $P(0, -5)$ oleh setengah putaran terhadap $O$.
        Apa aturan umum untuk $(x, y)$?
  \item Sekarang pusatnya $C(2, 1)$. Hitunglah bayangan $M(5, 3)$,
        lalu periksalah resep umumnya: setiap koordinat bayangannya
        adalah \emph{dua kali koordinat pusatnya dikurangi koordinat
        titiknya}.
  \item Kartu remi dirancang agar terbaca sama ketika dibalik ---
        sebuah \omterm{def:g7:central:def}{simetri pusat}, jadi tidak ada pemain yang melihat
        kartunya ``keliru''. Jelaskan mengapa, pada kartu dengan
        lambang berjumlah \emph{\omterm{def:g3:numbers:evenodd}{ganjil}}, satu lambang harus terletak
        tepat di pusat kartunya.
  \item Dua setengah putaran berturut-turut: ambillah pusat
        $O_1(2, 0)$ dan $O_2(5, 0)$. Kirimlah titik $A(1, 1)$ lewat
        setengah putaran terhadap $O_1$, lalu kirimlah bayangannya
        lewat setengah putaran terhadap $O_2$. Bandingkan awal dan
        akhirnya. Ulangi dengan $B(0, 3)$. Transformasi tunggal yang
        sederhana apa yang setara dengan kedua setengah putaran itu?
  \item Ukurlah geseran pertanyaan 4 terhadap jarak $O_1 O_2$, lalu
        nyatakan penemuanmu. (Bandingkan dengan dua cermin \omterm{def:g6:lines:perp}{sejajar}
        pada \cref{pb:g6:symmetry:1}: gejala yang sama, pemain yang
        baru.)
\end{enumerate}

\medskip
\noindent\textbf{Bagian II --- Kue dan garis berat.}
\begin{enumerate}[resume]
  \item Misalkan $O$ pusat sebuah \omterm{def:g7:central:parallelogram}{jajargenjang}. Jelaskan mengapa
        setiap garis melalui $O$ memotong tepinya di dua titik yang
        \omterm{def:g7:central:def}{simetris} terhadap $O$, dan mengapa garis itu memotong
        \omterm{def:g7:central:parallelogram}{jajargenjangnya} menjadi dua bagian yang persis merupakan
        salinan setengah putaran satu sama lain --- jadi luasnya sama.
  \item Teorema tukang roti: kue berbentuk persegi panjang terbagi
        dengan sangat adil oleh potongan lurus \emph{apa pun} yang
        melalui pusatnya. Lebih baik lagi: kue persegi panjang itu
        punya lubang persegi panjang dengan ukuran apa pun, letak apa
        pun, bahkan miring. Uraikan satu potongan lurus yang memberi dua
        bagian dengan kue yang sama banyak --- lalu berilah alasannya
        dengan pertanyaan 6.
  \item Tiga \omterm{def:g5:solids:def}{titik sudut} \omterm{def:g7:central:parallelogram}{jajargenjang} $ABCD$ adalah $A(0, 0)$,
        $B(6, 1)$ dan $C(8, 5)$. Hitunglah koordinat $D$
        dan koordinat pusatnya $O$.
  \item \Cref{exo:g7:central:9} membangun, dari segitiga $ABC$
        dan titik tengah $O$ dari $[BC]$, \omterm{def:g7:central:parallelogram}{jajargenjang}
        $ABA'C$ dengan $A'$ titik \omterm{def:g7:central:def}{simetris} $A$ terhadap $O$.
        Pakailah itu, bersama ketaksamaan segitiga
        (\cref{thm:g7:triangles:inequality}) pada segitiga
        $ABA'$, untuk membuktikan \emph{batas garis berat}: garis
        berat $[AO]$ memenuhi
        \[
        AO < \frac{AB + AC}{2} .
        \]
        (Perhatikan bahwa $AA' = 2\,AO$ dan $BA' = AC$.)
  \item Pada segitiga dengan $AB = 5$ cm dan $AC = 7$ cm, di antara
        dua nilai berapa panjang garis berat dari $A$ harus
        terletak? (Pakailah kedua arah ketaksamaan segitiga pada
        segitiga $ABA'$.)
\end{enumerate}

\medskip
\noindent\textbf{Bagian III --- Berapa banyak pusat yang dapat dipunyai
sebuah bangun?}
\begin{enumerate}[resume]
  \item Untuk tiap bangun, katakan apakah ia punya \omterm{def:g7:central:center}{pusat simetri},
        dan di mana: sebuah \omterm{def:g6:lines:objects}{ruas garis}; garis penuh; \omterm{def:g6:shapes:triangles}{segitiga sama
        sisi}; persegi; lingkaran; huruf S dan Z.
  \item Buktikan pernyataanmu untuk \omterm{def:g6:shapes:triangles}{segitiga sama sisi}: setengah
        putaran yang mengirim segitiga itu ke dirinya sendiri akan
        memasangkan ketiga \omterm{def:g5:solids:def}{titik sudutnya} --- tetapi tiga itu \omterm{def:g3:numbers:evenodd}{ganjil},
        jadi ada \omterm{def:g5:solids:def}{titik sudut} yang terkirim ke dirinya sendiri. \omterm{def:g5:solids:def}{Titik
        sudut} itu terpaksa menjadi apa, dan mengapa itu mustahil?
  \item Perumumlah pertanyaan 12: jelaskan mengapa \emph{tidak ada}
        \omterm{def:g4:geometry:polygon}{segi banyak} dengan \omterm{def:g5:solids:def}{titik sudut} berjumlah \omterm{def:g3:numbers:evenodd}{ganjil} yang punya
        \omterm{def:g7:central:center}{pusat simetri}.
  \item Mutiaranya: andaikan sebuah bangun punya \emph{dua} \omterm{def:g7:central:center}{pusat
        simetri} yang berbeda, $O_1$ dan $O_2$. Lakukan kedua setengah
        putaran itu berturut-turut lalu pakailah penemuan pada
        pertanyaan 5: gerak apa yang harus mengirim bangunnya tepat
        ke dirinya sendiri? Jelaskan mengapa bangun yang muat pada
        selembar kertas tidak dapat bertahan terhadap itu, lalu
        simpulkan.
  \item Kesimpulan pertanyaan 14 meninggalkan celah: bangun yang
        \emph{tak terbatas}. Periksalah pada pita hias tak berhingga
        --- misalnya, jalur jejak kaki tanpa ujung: kiri, kanan,
        kiri, kanan\,\dots\ --- bahwa ia punya (sedikitnya)
        dua \omterm{def:g7:central:center}{pusat simetri} yang berbeda, lalu kenalilah translasi
        yang sudah diramalkan pertanyaan 14.

\end{enumerate}
\end{problem}
